\documentclass[11pt,a4paper]{article}
\usepackage[T1]{fontenc}
\usepackage[utf8]{inputenc}
\usepackage{lmodern,microtype}
\usepackage[margin=27mm,headheight=14pt]{geometry}
\usepackage{amsmath,amssymb,amsthm,mathtools}
\usepackage{booktabs,longtable,array,enumitem}
\usepackage{xcolor,fancyhdr}
\usepackage[colorlinks=true,linkcolor=blue!45!black,urlcolor=blue!45!black,citecolor=blue!45!black]{hyperref}
\usepackage{listings}
\lstset{basicstyle=\small\ttfamily,breaklines=true,columns=fullflexible,frame=single,rulecolor=\color{black!15},xleftmargin=3pt,xrightmargin=3pt}
\setlength{\parskip}{4pt}
\setlength{\parindent}{0pt}
\setlength{\emergencystretch}{2em}
\setlist{itemsep=3pt,topsep=5pt}
\pagestyle{fancy}
\fancyhf{}
\fancyhead[L]{\small Constructive Diophantine approximation in iPRA}
\fancyhead[R]{\small Hurwitz and Khinchin}
\fancyfoot[C]{\thepage}
\newtheorem{theorem}{Theorem}[section]
\newtheorem{lemma}[theorem]{Lemma}
\newtheorem{proposition}[theorem]{Proposition}
\theoremstyle{definition}
\newtheorem{definition}[theorem]{Definition}
\newcommand{\N}{\mathbb N}
\newcommand{\Q}{\mathbb Q}
\newcommand{\Z}{\mathbb Z}
\newcommand{\ipra}{\mathsf{iPRA}}
\newcommand{\PR}{\mathsf{PR}}
\newcommand{\len}{\operatorname{len}}
\newcommand{\ceil}[1]{\left\lceil #1\right\rceil}
\newcommand{\repo}{\texttt{ProveIt}}
\title{\vspace{-12mm}\textbf{Constructive Hurwitz and Khinchin theorems\\in intuitionistic primitive recursive arithmetic}\\[5pt]\large Finite witnesses, rational interval covers, and a Lean formalization route}
\author{Research report prepared for Vladimir Reshetnikov}
\date{15 November 2026}
\begin{document}
\maketitle
\begin{abstract}
We give fully developed finite proofs of constructive versions of two
Diophantine approximation theorems.  Hurwitz's sharp coefficient
$1/\sqrt5$ is retained: from a primitive recursive rational name and an
explicit rational-apartness modulus, a primitive recursive procedure
returns a reduced approximation beyond any prescribed denominator,
together with a strict rational certificate.  For Khinchin's theorem,
convergence with a supplied tail modulus gives explicit small tail covers.
Divergence with a supplied unit-tail modulus gives a finite denominator
cutoff whose approximation intervals cover all but any prescribed dyadic
tolerance of $[0,1]$.  The divergence proof uses elementary rational-point
counting, a finite second-moment inequality, and a fixed number of rounds
of filling residual gaps.  An explicit null-test interpretation preserves
positive approximation witnesses without a classical quantifier-negation
step.  We explain why all finite assertions and algorithms admit proofs
using quantifier-free induction in $\ipra$, and give a module-by-module
route to Lean and the deeply embedded \repo{} iPRA calculus.
\end{abstract}

\paragraph{Status of this package.}
The mathematical arguments are supplied below.  The Python artifacts are
exact-arithmetic implementations and finite checks of the certificate
interfaces and principal finite estimates.  The Lean artifact is an
implementation scaffold and an explicit list of proof obligations; it is
\emph{not} a completed Lean formalization of either theorem, nor an
embedded iPRA proof certificate.  No Lean compiler was available in the
execution environment, so the scaffold has not been compiled.  These
boundaries are part of the deliverable, not assumptions used in the proofs.

\clearpage
\tableofcontents
\clearpage
\section{The precise constructive setting}
\label{sec:foundations}

\subsection{The statements being proved}

The linked Hurwitz theorem concerns infinitely many reduced fractions
$p/q$ satisfying $|x-p/q|<1/(\sqrt5 q^2)$ for each irrational $x$.
The linked Khinchin theorem concerns almost-everywhere approximation by
$|x-p/q|<\psi(q)/q$ and the series $\sum_q\psi(q)$.
Their literal classical formulations quantify over real numbers, infinite
sequences, and Lebesgue measurable sets.  First-order $\ipra$ has one
numerical sort, primitive recursive function symbols, intuitionistic
first-order logic with equality, and induction only for quantifier-free
formulas.  We therefore specify the finite arithmetic content of the
constructive versions before using real-number language.

The two central output assertions are as follows.
\begin{enumerate}
\item \textbf{Hurwitz witness.} For a fixed primitive recursive rational
name $x_s$ with error at most $2^{-s}$, a fixed primitive recursive
rational-apartness modulus, and every natural $B$, compute $p,q,s$ with
$q>B$, $\gcd(p,q)=1$, and
\[
 5q^4\bigl(|x_s-p/q|+2^{-s}\bigr)^2<1.
\]
This is a decidable inequality between rational numbers.  Its proof gives
the strict classical coefficient, rather than a coefficient enlarged by
an arbitrary epsilon.
\item \textbf{Khinchin finite coverage.} For fixed nonnegative rational
primitive recursive $\psi$ and an appropriate primitive recursive modulus,
compute finite rational interval lists.  In the divergent case the
output $U(N,m)$ is labelled by reduced fractions with denominators
between $N$ and a computed $M(N,m)$ and satisfies
\[
 U(N,m)\subseteq\bigcup_{q=N}^{M(N,m)}
       \bigcup_{p=0}^q (p/q-\psi(q)/q,p/q+\psi(q)/q),
 \qquad |[0,1]\setminus U(N,m)|<2^{-m}.
\]
In the convergent case every finite portion of a specified tail has
length at most the specified tolerance.  The measure symbols here mean
finite rational length computed by endpoint sorting.
\end{enumerate}

The divergence proof needs only that $\psi$ be nonincreasing.  The linked
encyclopedia page states the stronger older hypothesis that $q\psi(q)$
is nonincreasing, in the same normalization $\psi(q)/q$.  That stronger
hypothesis implies ours: for $q\le r$, $r\psi(r)\le q\psi(q)$ implies
$\psi(r)\le(q/r)\psi(q)\le\psi(q)$.  The present proof consequently
covers the linked formulation as well as the usual weaker monotonicity
assumption.  The conversion to the alternative radius notation
$\Psi(q)=\psi(q)/q$ is stated separately below.

\subsection{Data, hypotheses, and theorem schemata}

``Primitive recursive'' always means obtained from zero, successor and
projections by composition and primitive recursion.  Natural numbers
code signed integers, fractions with positive denominator, finite lists,
and finite machine states.  Numerator signs are stored explicitly.
Rational arithmetic is implemented by integer arithmetic and comparison
by cross multiplication.  Reduction of fractions is optional except where
coprimality is required; gcd can be computed by the bounded search through
divisors up to the larger input.  List decoding and sorting can be performed
by bounded loops on the code and its recorded length.  Efficiency is not
needed for any foundational assertion here.

An input function such as $\psi$, $x_s$, or a modulus is a \emph{fixed}
primitive recursive definition, represented by its own function symbol.
The statements are schemata, one for every such fixed collection of
definitions.  Their hypotheses need not already be theorems of $\ipra$:
the formal statement is the implication from the displayed coherence,
monotonicity, separation, or mass assumptions.  For a particular example,
one additionally proves those assumptions.  Alternatively all queried
values and their finite verifications may be supplied as finite tables.
There is no first-order quantification over arbitrary real functions and
no supposed primitive recursive universal evaluator for all primitive
recursive program codes.

A name $s\mapsto x_s\in\Q$ may be specified arithmetically by
\[
 |x_s-x_t|\le2^{-s}+2^{-t}\qquad(s,t\in\N).
\]
The notation $|x-x_s|\le2^{-s}$ expresses its ordinary interpretation;
the internal arguments use rational enclosures and these coherence
inequalities.  A positive separation modulus supplies quantitative
apartness from every rational with bounded denominator.  Mere negative
irrationality, expressed as the absence of an equality, is not being
silently promoted to that modulus.

Similarly we do not decide whether an arbitrary nonnegative series
converges or diverges.  We prove separate implications from a supplied
tail bound or a supplied divergence witness.  No procedure extracting a
primitive recursive modulus from an unqualified classical convergence or
divergence assertion is part of either theorem.  Section~\ref{sec:extensions}
extends the metric argument to real-valued approximation functions with
primitive recursive names and witnessed tail information.

\subsection{Bounded predicates and quantifier-free induction}

Every bounded predicate on the finite data used here has a primitive
recursive Boolean characteristic function.  For a Boolean predicate
$b(i)$, conjunction over $i<n$ is the finite product of the bits, and
existence is a bounded search with a failure flag.  Nesting these
constructions handles finitely nested bounded quantifiers.

There is a relevant proof-theoretic detail.  We do not apply the iPRA
induction rule directly to a formula containing a bounded quantifier.
We first replace its bounded tests by the corresponding characteristic
terms, prove their specifications, and induct on the resulting equation.
For example the equations
\[
 a(0)=1,\qquad a(n+1)=a(n)b(n)
\]
are accompanied by a bounded first-failure function.  Quantifier-free
induction proves that a failure flag identifies an index $i<n$ with
$b(i)=0$, and that $a(n)=1$ forces $b(i)=1$ for each given $i<n$.
The converse follows from the failure witness and the decidable bit
case distinction.  This proves the bounded-universal specification
without induction on a quantified formula.  The same method handles
finite maxima, bounded searches, and endpoint membership tests.

More generally, let $S(0,z)=G(z)$ and
$S(n+1,z)=F(n,S(n,z),z)$ be a fixed primitive recursive state iteration.
If $V(n,s,z)$ is a primitive recursive Boolean invariant and its initial
and preservation equations have been proved, iPRA's quantifier-free
induction proves $V(n,S(n,z),z)=1$.  Universal assumptions about the input
names or moduli may remain in the context and be instantiated at each
step; they are not induction formulas.  The invariants used below are
listed explicitly in Section~\ref{sec:formalization}.  Thus the reasoning
is a finite arithmetic proof with a specified induction discipline, not
an appeal to the rule that every terminating algorithm has iPRA-provable
correctness.  That rule would be unjustified.

\subsection{Finite length and finite second moments}

For rational open intervals, retain their endpoints and their open/closed
flags when membership matters.  Sort the finite list of all endpoints.
Between successive distinct endpoints, membership in each input interval
is constant and can be evaluated at the rational midpoint.  The length
of a union is the sum of the lengths of exactly those cells whose
membership bit is one.  Single endpoints have length zero, but a singleton
in a complement must still be retained for later open-cover statements.
This procedure proves finite additivity on disjoint cells and finite
subadditivity by the inequality between membership bits.

If intervals only touch, their open union omits the common endpoint.
It is permissible to merge strict overlaps; merging touching intervals
would lose that singleton.  The algorithms and proof below retain it.
The two endpoints of $[0,1]$ may safely be included as extra singleton
gaps, even if a relatively open interval already covers them.

For a nonnegative integer-valued function $F$ constant on these cells,
the notation $\int F$ abbreviates a finite rational sum
$\sum_i w_iF_i$.  In particular the only Cauchy--Schwarz inequality needed
in the metric proof follows from
\[
 \left(\sum_i w_i\right)\left(\sum_i w_i v_i^2\right)
   -\left(\sum_i w_i v_i\right)^2
 =\sum_{i<j}w_iw_j(v_i-v_j)^2\ge0.
\]
This is a polynomial identity and a sum of nonnegative rational terms.
It uses no integration theory, completeness theorem, or countable choice.
All appearances of measure inside the main proof are of this form.

\subsection{Relation to the classical literature}

The sharp Hurwitz theorem and its three-convergent proof are classical;
an accessible author-written account is Filaseta's Theorem~25
\cite{filaseta}.  The organization of the metric proof belongs to the
local-distribution and overlap methods underlying local ubiquity
\cite{bv,bdv}.  We prove the counting, overlap, and amplification estimates
in full, rather than importing the general measure-theoretic theorem.
The specific constants are chosen to simplify verification.  We make no
claim that these constants are optimal, or that the classical mathematical
ideas are new.

% Standalone section body. Requires amsmath, amssymb, hyperref only.
\section{Hurwitz approximation with a finite witness}
\label{sec:constructive-hurwitz}

The classical constant is retained exactly.  The constructive change is to
replace an unqualified irrational real by a fast rational name together with a
quantitative certificate that it stays away from rationals.  The output is a
rationally checkable strict inequality, not a double-negated existence claim.

\subsection{The input and the exact statement}

Fix primitive recursive functions \(s\mapsto x_s\in\mathbb Q\) and
\(Q\mapsto I(Q)\in\mathbb N\).  In the readable notation below, they name a
real \(x\) and satisfy
\begin{equation}\label{eq:hw-name}
 |x-x_s|\le 2^{-s},\qquad
 |x-p/q|\ge 2^{-I(Q)}
 \quad(p\in\mathbb Z,\ 1\le q\le Q).
\end{equation}
We take \(I\) nondecreasing, replacing it by
\(\max_{1\le j\le Q}I(j)\) when necessary.  This replacement is primitive
recursive.  ``Real'' here is explanatory notation: the first-order formulation
uses rational Cauchy inequalities and, for example, the rational separation
conditions
\[
 |x_s-p/q|+2^{-s}\ge2^{-I(Q)}
 \qquad(1\le q\le Q)
\]
for every indicated index.  Only finitely many instances are used for each
output.  There is no sort of arbitrary real numbers in iPRA.

\paragraph{Constructive Hurwitz theorem.}
For every \(B\in\mathbb N\), a primitive recursive computation from the fixed
functions \(x_s,I\) produces \(p\in\mathbb Z\), \(q,s\in\mathbb N\) such that
\begin{equation}\label{eq:hw-certificate}
 B<q,\qquad \gcd(p,q)=1,\qquad
 5q^4\bigl(|x_s-p/q|+2^{-s}\bigr)^2<1.
\end{equation}
In particular,
\[
 |x-p/q|<\frac1{\sqrt5\,q^2}.
\]
Thus the theorem provides approximants beyond every prescribed denominator,
and therefore an explicitly generated sequence of distinct approximants.
One denominator bound is
\begin{equation}\label{eq:hw-R}
 R_0=1,\qquad R_{j+1}=\max\{R_j,2^{I(R_j)}\},\qquad
 q\le R_{B+3}.
\end{equation}
No useful efficiency is claimed for this bound; its purpose is to make the
primitive recursive character explicit.

As usual for a first-order theory, this is a theorem schema for fixed primitive
recursive definitions of the name and modulus.  It is not an assertion that
iPRA contains a total universal evaluator for all primitive recursive programs.
An arbitrary computable irrational need not supply primitive recursive data
of the required kind.  Dropping the quantitative separation hypothesis would
change the constructive content of the statement.

\subsection{Finite continued fractions and their identities}

Take integers \(a_0\) and \(a_j\ge1\) for \(j\ge1\), and put
\[
 p_{-1}=1,\quad p_0=a_0,\qquad q_{-1}=0,\quad q_0=1,
\]
\[
 p_j=a_jp_{j-1}+p_{j-2},\qquad
 q_j=a_jq_{j-1}+q_{j-2}\quad(j\ge1),
\]
where the recurrence at \(j=1\) uses the displayed initial values in the
usual way (equivalently start with \(p_{-2}=0,q_{-2}=1\)).
Induction on the finite recurrence gives
\begin{equation}\label{eq:hw-det}
 p_jq_{j-1}-p_{j-1}q_j=(-1)^{j-1}\quad(j\ge1).
\end{equation}
Consequently \(p_j,q_j\) are coprime, and so are \(q_j,q_{j-1}\).
Also \(q_1\ge1\), \(q_2\ge2\), and
\begin{equation}\label{eq:hw-growth}
 q_j\ge j\quad(j\ge1),\qquad q_{j+1}\ge q_j+q_{j-1}.
\end{equation}
For the determinant identity, substituting the recurrence changes its sign;
its first value is \(1\).  For the growth bound, \(q_2\ge2\) starts the
induction and the recurrence adds at least one at each later step.

Suppose these are successive digits of the continued fraction of \(x\).
The complete quotient \(t_{j+1}>1\) satisfies
\[
 x=\frac{p_jt_{j+1}+p_{j-1}}{q_jt_{j+1}+q_{j-1}},\qquad
 t_{j+1}>a_{j+1}.
\]
This identity follows by substituting \(t_j=a_j+1/t_{j+1}\) repeatedly;
only a finite number of substitutions is needed.  Subtracting \(p_j/q_j\)
and using \eqref{eq:hw-det} gives
\begin{equation}\label{eq:hw-error}
 \left|x-\frac{p_j}{q_j}\right|
 =\frac1{q_j(q_jt_{j+1}+q_{j-1})}
 <\frac1{q_jq_{j+1}}.
\end{equation}
The sign alternates, so \(x\) lies strictly between consecutive convergents.
For adjacent denominators \(u,v\) and their corresponding numerators \(r,t\),
\begin{equation}\label{eq:hw-bracket}
 \left|x-\frac r u\right|+
 \left|x-\frac t v\right|=\frac1{uv}.
\end{equation}
These are identities about finitely many rational fractional-linear
transformations, with a positive remainder.  They do not require a theorem
about the existence of an infinite continued fraction.

Combining the separation assumption with \eqref{eq:hw-error} yields
\[
 2^{-I(q_j)}\le |x-p_j/q_j|<\frac1{q_jq_{j+1}},
 \qquad q_{j+1}<2^{I(q_j)}.
\]
This proves \eqref{eq:hw-R} by finite induction once the necessary prefix
has been constructed.

\subsection{A rational replacement for the three-convergent argument}

The next lemma is the key to retaining the sharp constant without deciding
comparisons between arbitrary reals.

\paragraph{Denominator lemma.}
Let \(1\le u\le v\), \(w\ge u+v\), and \(\gcd(u,v)=1\).
At least one of the two rational numbers
\[
 S_0=\frac{u^2+v^2}{uv},\qquad
 S_1=\frac{v^2+w^2}{vw}
\]
has square strictly greater than \(5\).  A suitable one is selected by an
integer comparison.

\paragraph{Proof.}
Set \(h=v^2-uv-u^2\).  It cannot vanish.  Indeed, if it did then
\(u\mid v^2\); coprimality implies \(u=1\), and \(v^2=v+1\) has no
positive integral solution: it fails at \(v=1\), while \(v\ge2\) implies
\(v^2\ge v+2\).
If \(h>0\), then
\[
 (u^2+v^2)^2-5u^2v^2
 =(v^2-uv-u^2)(v^2+uv-u^2)>0,
\]
since \(v^2+uv-u^2\ge uv>0\).  Choose \(S_0\).
If \(h<0\), then
\[
 w^2-vw-v^2\ge(u+v)^2-v(u+v)-v^2=-h>0.
\]
The inequality follows either by expansion or from the exact nonnegative
remainder \((w-u-v)(w+u)\).  Applying the preceding factorization to
\(v,w\) gives \(S_1^2>5\).  Choose \(S_1\).
All comparisons in this proof are integer comparisons. \hfill\(\square\)

Apply the lemma to
\[
 (u,v,w)=(q_{B+1},q_{B+2},q_{B+3}).
\]
Every denominator in this triple exceeds \(B\).  Denote the selected pair
again by \(u,v\), and its numerators by \(r,t\).  Put
\begin{equation}\label{eq:hw-margin}
 C=\frac{uv}{u^2+v^2},\qquad d=1-5C^2>0.
\end{equation}
In particular \(0<C\le1/2\), \(0<d<1\), and \(C<1/\sqrt5\).
The familiar classical proof now says that one of the two normalized
errors is at most \(C\).  We instead give a finite selection procedure
that includes a margin certificate.

Let \(Q=\max(u,v)\) and choose a precision \(s\) satisfying
\begin{equation}\label{eq:hw-precision}
 Q^2 2^{-s}\le d/100.
\end{equation}
This need not be found by an unbounded search.  Since
\[
 d=\frac{(u^2+v^2)^2-5u^2v^2}{(u^2+v^2)^2}
 \ge\frac1{(u^2+v^2)^2},
\]
the explicit integer
\[
 s=100Q^2(u^2+v^2)^2
\]
suffices, using \(2^s\ge s\).
Compute the two rational numbers
\[
 A=u^2|x_s-r/u|,\qquad D=v^2|x_s-t/v|
\]
and choose the smaller (with a fixed tie rule).  The triangle inequality
and \eqref{eq:hw-bracket} give
\[
 \frac A{u^2}+\frac D{v^2}
 \le\frac1{uv}+2\cdot2^{-s}.
\]
Therefore their smaller value is at most
\[
 \frac{1/(uv)+2\cdot2^{-s}}{1/u^2+1/v^2}
 =C+2\cdot2^{-s}\frac{u^2v^2}{u^2+v^2}
 \le C+Q^2 2^{-s}.
\]
For the selected numerator and denominator \(p,q\), it follows that
\begin{equation}\label{eq:hw-cert-bound}
 q^2\bigl(|x_s-p/q|+2^{-s}\bigr)
 \le C+2Q^2 2^{-s}\le C+d/50.
\end{equation}
Finally,
\[
 5(C+d/50)^2
 =1-d+Cd/5+d^2/500
 \le1-d+d/10+d/500<1.
\]
Together with \eqref{eq:hw-cert-bound}, this is exactly the rational
certificate \eqref{eq:hw-certificate}.  In particular the proof produces
one of every three indicated convergents and does not use excluded middle
for the real errors.  The constant \(\sqrt5\) disappears entirely from
the certificate checker.

\subsection{Why the continued-fraction prefix can be obtained primitively recursively}

For completeness we give an explicit bounded-precision extraction method.
This is the point at which qualitative irrationality would be insufficient.
Let \(X=\lceil|x_0|\rceil+2\), so \(|x|<X\).
To find \(a_0=\lfloor x\rfloor\), take
\(s_0=I(1)+3\) and compute \(a_0=\lfloor x_{s_0}\rfloor\).
No integer lies within \(2^{-I(1)}\) of \(x\), whereas
\(|x-x_{s_0}|\le2^{-I(1)}/8\).  Thus these two points have the same
integer part.  This also certifies \(0<x-a_0<1\).

Assume the prefix through \(a_{n-1}\) has been obtained, where \(n\ge1\),
and write
\[
 P=p_{n-1},\ Q=q_{n-1},\quad P'=p_{n-2},\ Q'=q_{n-2}.
\]
Here \(Q>0\); for \(n=1\) we use \(P'=1,Q'=0\).
The next complete quotient is
\begin{equation}\label{eq:hw-mobius}
 t_n=\frac{P'-Q'x}{Qx-P}>1.
\end{equation}
Its denominator is nonzero, with the known lower bound
\[
 |Qx-P|\ge D_0:=Q2^{-I(Q)}>0.
\]
Set
\[
 E=\max(1,QX+|P|),\qquad
 U=\left\lceil\frac{|P'|+Q'X}{D_0}\right\rceil+2.
\]
Then \(|Qx-P|\le E\) and \(1<t_n<U-1\).
For every integer \(k\) with \(0\le k\le U\),
\[
 t_n-k=\frac{(P'+kP)-(Q'+kQ)x}{Qx-P}.
\]
Put \(H=Q'+UQ\) and \(L=2^{-I(H)}/E\).
If \(Q'+kQ>0\), irrational separation shows that the absolute value
of the numerator is at least \(2^{-I(H)}\).
The only possible zero-denominator coefficient is \(Q'=k=0\), when
\(|P'|=1\) by the determinant identity, giving the same lower bound.
Consequently
\begin{equation}\label{eq:hw-floor-gap}
 |t_n-k|\ge L>0\qquad(0\le k\le U).
\end{equation}

Choose a rational precision bound
\[
 \delta=\min\{D_0/(4Q),\ LD_0^2/8\}>0,
\]
and choose \(s\) with \(2^{-s}\le\delta\) by a bounded rational computation.
For instance, if \(\delta=a/b>0\) in integer form, \(s=b+1\) suffices;
no reduction of the fraction is necessary.  Put \(y=x_s\).
The first bound gives \(|Qy-P|\ge3D_0/4\).  The determinant identity gives
the exact difference formula
\[
 \left|
 \frac{P'-Q'y}{Qy-P}-\frac{P'-Q'x}{Qx-P}
 \right|
 =\frac{|x-y|}{|Qy-P|\,|Qx-P|}
 \le\frac{4\delta}{3D_0^2}\le L/6.
\]
Here $L\le1$, and $1<t_n<U-1$ ensures that only the integers in
$[0,U]$ can matter.  By \eqref{eq:hw-floor-gap}, the two quotients have the same integer part.
Thus exact rational division and flooring of the quotient at \(y\) gives
\(a_n\), and the associated inequalities certify
\(a_n<t_n<a_n+1\).  Set \(t_{n+1}=1/(t_n-a_n)\) in the explanatory
notation and update the integer convergents by their recurrences.

Iterating this finite procedure through \(a_{B+3}\) supplies the prefix
used above.  Every operation is integer or rational arithmetic, a bounded
maximum, fixed-length primitive recursion, or evaluation of the fixed
primitive recursive input functions.  Each claimed real inequality has a
strict rational margin and can be replaced by inequalities between sufficiently
accurate rational names.  The loop invariants consist of determinant,
positivity, and rational enclosure certificates; they are decidable finite
conditions.  This is the promised path to a quantifier-free induction proof.

An alternative interface accepts a primitive recursive sequence of positive
continued-fraction digits directly.  Its convergents give a primitive recursive
Cauchy name (the denominators grow at least as Fibonacci numbers).  A
quantitative irrationality witness also follows: for \(1\le q\le B\), choose
an explicitly indexed convergent with \(q_n>B\) and \(q_{n+1}\ge2B\).
Then, unless the rational equals that convergent (which is impossible in
lowest terms),
\[
 |x-p/q|\ge\frac1{Bq_n}-\frac1{q_nq_{n+1}}
 \ge\frac1{2Bq_n}.
\]
A later index such as \(n=2B+2\) suffices by \eqref{eq:hw-growth}.
This explains the equivalence of the two useful input interfaces without
postulating an unbounded continued-fraction oracle.
For example, $I(B)=2Bq_{2B+2}$, for $B\ge1$, is an explicit sufficient
modulus for the digit interface, since $2^r\ge r$; it may be replaced
by its bounded running maximum.

\subsection{An entirely finite certificate interface}

The use of $x$ above is eliminable at the main arithmetic step.
Supply the triple of integer convergents and a rational interval
$[\ell,h]$ contained in both adjacent-convergent brackets.  Its midpoint
$y$ and half-width $e$ satisfy $|z-y|\le e$ for every $z$ in that interval.
For the pair selected by the denominator lemma, require the rational
test $Q^2e\le d/100$.  The same weighted-error calculation, now entirely
over rational endpoints, selects $p/q$ and proves
\[
 5q^4\max\{|\ell-p/q|,|h-p/q|\}^2<1.
\]
The convexity inequality for absolute value on an interval follows by
the three rational cases $p/q\le\ell$, $\ell<p/q<h$, and $h\le p/q$.
Consequently this is a finite rational theorem with no real parameters.

In the name interface, refine the final precision further so that
$2^{-s}\le2^{-I(Q_*)-3}$, where $Q_*$ is the largest denominator of
the triple.  Each convergent endpoint stays at least $2^{-I(Q_*)}$
from the named value.  Thus $[x_s-2^{-s},x_s+2^{-s}]$ lies inside both
adjacent-convergent brackets, and the enclosure checks are finite rational
checks.  The added precision is a primitive recursive maximum of the
two prescribed bounds and does not change the denominator bound.

The same enclosure discipline applies during digit extraction.  With
the displayed choice of $\delta$, every rational $z$ in
$[y-2^{-s},y+2^{-s}]$ has $|z-x|\le2\delta$ in the explanatory
notation, hence $|Qz-P|\ge D_0/2$.  The fractional-linear value changes
by at most $4\delta/D_0^2\le L/2$, so the rational image interval
crosses no integer.  Its two endpoint floors certify the next digit.
Coherence of the name supplies compatible later refinements.  This
gives the concrete finite-enclosure invariant to be implemented in
the object-theory formalization.

\subsection{Optimality, with an explicit finite denominator bound}

Let \(\phi=(1+\sqrt5)/2\) and \(\phi'=(1-\sqrt5)/2\).
These numbers have primitive recursive rational names by bisection of
\(z^2-5\) on \([2,3]\), using a prescribed finite number of bisections.
For every \(p\in\mathbb Z,q\ge1\),
\begin{equation}\label{eq:hw-norm}
 N=p^2-pq-q^2\ne0,\qquad
 1\le|N|=q^2|p/q-\phi|\,|p/q-\phi'|.
\end{equation}
To verify that \(N\ne0\), divide \(p,q\) by their greatest common divisor.
If \(N=0\), coprimality would force the reduced denominator to be one,
and an integer \(z\) would satisfy \(z(z-1)=1\).  The product is zero
at \(z=0,1\), and at least two for every other integer, a contradiction.
This proof uses only integer arithmetic.

The same identity supplies the required irrationality modulus.  If
\(0\le p/q\le3\), then \(|p/q-\phi'|<4\), whence
\(|\phi-p/q|>1/(4q^2)\).  Outside that rationally decidable range the
distance is greater than one.  Hence
\[
 |\phi-p/q|\ge\frac1{4Q^2}\ge2^{-4Q^2}
 \qquad(1\le q\le Q),
\]
so \(I(Q)=4Q^2\) is a valid elementary modulus.

Now let \(c>0\) be rational with \(5c^2<1\), and put
\(d=1-5c^2>0\).  Suppose
\[
 |\phi-p/q|<c/q^2.
\]
Since \(|\phi-\phi'|=\sqrt5\), \eqref{eq:hw-norm} gives
\[
 1<c\sqrt5+c^2/q^2.
\]
But
\[
 1-c\sqrt5=\frac{1-5c^2}{1+c\sqrt5}>d/2.
\]
Therefore
\begin{equation}\label{eq:hw-optimal-bound}
 q^2<\frac{2c^2}{1-5c^2}.
\end{equation}
Only finitely many denominators are possible; for each denominator the
approximation inequality also bounds the numerator.  Every proposed rational
coefficient strictly smaller than \(1/\sqrt5\) therefore fails to give
approximants beyond every denominator for this explicitly presented irrational.
The argument proves optimality constructively, with the explicit bound
\eqref{eq:hw-optimal-bound}, rather than merely by a limiting statement about
Fibonacci ratios.

\subsection{Lean implementation plan for this part}

The arithmetic core should be formalized over \(\mathbb Z\) and a rational
number representation, not by importing a classical real-number Hurwitz
theorem.  A useful dependency order is:
\begin{enumerate}
\item finite signed-integer and rational arithmetic, powers of two, bounded
      maxima, and exact rational comparison;
\item convergent recurrences and the determinant, coprimality, and growth
      identities;
\item the denominator lemma, using its two displayed polynomial
      factorizations and the decidable sign of \(h\);
\item a rational bracketing lemma: a point in an interval of length \(1/(uv)\)
      and its \(\varepsilon\)-approximation satisfy the weighted-error bound;
\item the selector and rational certificate \eqref{eq:hw-certificate}, with
      \(C,d,s\) exactly as above;
\item bounded-precision prefix extraction from a name and a separation
      modulus, maintaining finite interval and determinant certificates;
\item the integer norm argument for the golden ratio and its bounded
      bisection name.
\end{enumerate}
The main internal statement can use the rational inequality
\(5q^4(|x_s-p/q|+2^{-s})^2<1\) directly.  A separate interpretation lemma
then reads this as the familiar real inequality.  Lean's ambient logic does
not by itself certify an iPRA derivation: the extraction functions need
primitive-recursion certificates, and their arithmetic proofs need to be
compiled into the chosen iPRA syntax or checked against its induction
restrictions.  The finite identities and invariants above make those
obligations explicit.

\paragraph{Classical reference.}
For the classical three-convergent theorem and the sharp constant, see
Michael Filaseta, \emph{Math 784: Algebraic Number Theory, Instructor's Notes},
Theorem 25, printed pages 20--21,
\url{https://people.math.sc.edu/filaseta/gradcourses/TheMath784Notes.pdf}.
The rational pair selector, explicit margins, certificate interface, and
bounded-precision extraction above give the constructive implementation
needed here; the classical statement alone would not supply it.

% Integration notes: finite constructive Khintchine theorem.
% This file has no preamble; the root manuscript may import or adapt it.
\section{A finite quantitative Khintchine theorem}
\label{sec:finite-khintchine}

We use the normalization
\[
 \left|x-\frac pq\right|<\frac{\psi(q)}q,
 \qquad q\geq1,
\]
so that the critical series is $\sum_q\psi(q)$.  In the normalization
$|x-p/q|<\Psi(q)$, put $\psi(q)=q\Psi(q)$: the monotonicity hypothesis is
that $q\Psi(q)$ is nonincreasing and the critical series is
$\sum_q q\Psi(q)$.  All approximation functions in the finite theorem
are nonnegative rational valued.  A primitive recursive approximation
function and primitive recursive modulus give a primitive recursive
construction; no uniform evaluator for all primitive recursive programs
is presupposed.

For a finite union of intervals with rational endpoints, write $|U|$ for
its exact rational length.  It is computed by sorting endpoints and adding
lengths of the covered elementary segments.  Single endpoints have length
zero.  Every measure argument below is an identity or inequality about
these finite rational sums.

\subsection{The hypotheses and conclusions}

For $q\geq1$, let $A_q$ be the relatively open subset of $[0,1]$ consisting
of the points satisfying the displayed approximation inequality for some
integer $p$.  Numerators $0\leq p\leq q$ suffice on $[0,1]$.

\paragraph{Convergent version.}
Suppose a function $C$ satisfies $C(s)\ge1$ and
\[
 \sum_{q=C(s)}^R\psi(q)\leq2^{-s}
 \quad\hbox{for every }s\hbox{ and every }R\geq C(s).
 \tag{K-C}
\]
Then, for every $m$ and $R\geq C(m+1)$,
\[
 \left|\bigcup_{q=C(m+1)}^R A_q\right|\leq2^{-m}.
 \tag{K1}
\]
Monotonicity of $\psi$ is not needed here.

\paragraph{Divergent version.}
Suppose $\psi$ is nonincreasing and a unit-tail modulus $D$ satisfies
\[
 D(s)\geq s,\qquad
 \sum_{q=s}^{D(s)}\psi(q)\geq1 \quad(s\geq1).
 \tag{K-D}
\]
There is a primitive recursive construction, relative to the fixed
primitive recursive functions $\psi,D$, that takes $N\geq1$ and $m\geq0$
and returns a finite list $U(N,m)$ of rational open intervals and an integer
$M(N,m)$ such that
\[
 U(N,m)\subseteq\bigcup_{q=N}^{M(N,m)}A_q,
 \qquad |[0,1]\setminus U(N,m)|<2^{-m}.
 \tag{K2}
\]
Every interval in the returned list is labelled by a reduced fraction
$p/q$ with $q>N$, and is contained in its approximation interval.  Thus the
construction also proves the version requiring coprime numerators and
denominators.

A usual primitive recursive initial-sum divergence rate $H$, with
$\sum_{q=1}^{H(t)}\psi(q)\geq t$, supplies (K-D): take
\[
 D(s)=\max\left(s,H\left(1+\left\lceil
       \sum_{q=1}^{s-1}\psi(q)\right\rceil\right)\right).
\]
All quantities in this formula are rational or natural numbers, and its
verification uses only nonnegativity and subtraction of finite sums.

\subsection{Convergence: the complete estimate}

Put $r=\psi(q)/q$.  The intervals centered at $p/q$, $1\leq p\leq q-1$,
have clipped length at most $2r$; the two intervals at $0$ and $1$ have
clipped length at most $r$ each.  Therefore
\[
 |A_q|\leq2qr=2\psi(q).
\]
Finite subadditivity and (K-C) immediately give (K1).  The same calculation
bounds every finite initial portion of the countable tail cover, which is
the constructive null-cover meaning of the convergence conclusion.  If
one insists on intervals open in all of $\mathbb R$ and counts their
unclipped lengths, the simpler bound $4\psi(q)$ applies; replace
$C(m+1)$ by $C(m+2)$.

\subsection{Enough separated reduced fractions in every interval}

Let $J=(a,b)\subseteq(0,1)$, allowing $a=0$ or $b=1$, and put $L=b-a>0$.
Its middle half is
\[
 J^*=(a+L/4,b-L/4),\qquad \lambda=|J^*|=L/2.
\]
For $Q=2^k$, $k\geq2$, let $\mathcal R_k(J)$ be the reduced fractions
$p/q\in J^*$ with $Q/4<q\leq Q$.  We prove that
\[
 \frac{LQ^2}{32}\leq\#\mathcal R_k(J)\leq LQ^2
 \tag{K3}
\]
whenever
\[
 k\geq K(J,N):=
 \max\left(4,512\lceil1/L\rceil,2+\lceil\log_2N\rceil\right).
 \tag{K4}
\]
Here $\lceil\log_2N\rceil$ is only notation for the least $j\leq N$
with $2^j\geq N$; it is a bounded search.  Its presence also ensures
$q>Q/4\geq N$ for the fractions used below.

Let $A(X)$ count all integer pairs $(p,q)$ with $1\leq q\leq X$ and
$p/q\in J^*$, without requiring coprimality.  An open interval of length
$t$ contains between $t-1$ and $t+1$ integers.  Consequently
\[
 \frac{\lambda X^2}{2}-X\leq A(X)
 \leq\frac{\lambda X^2}{2}+\left(\frac\lambda2+1\right)X.
 \tag{K5}
\]
Every noncoprime pair counted by $A(Q)$ is divisible by some integer
$d\in\{2,\ldots,Q\}$.  Dividing both coordinates by $d$ and applying
(K5) bounds the number of bad pairs by
\[
 \sum_{d=2}^Q A(\lfloor Q/d\rfloor)
 \leq\frac{\lambda Q^2}{2}\sum_{d=2}^Q\frac1{d^2}
      +\left(\frac\lambda2+1\right)Q\sum_{d=2}^Q\frac1d.
\]
No prime-number theorem, totient asymptotic, or inversion formula occurs
in this estimate.  The required elementary inequalities are
\[
 \sum_{d=2}^Q\frac1{d^2}\leq\frac14+
       \sum_{d=3}^Q\frac1{d(d-1)}\leq\frac34,
 \qquad
 \sum_{d=2}^{2^k}\frac1d\leq k.
 \tag{K6}
\]
The first telescopes, and the second splits the sum into the $k$ blocks
$(2^{j-1},2^j]$, each of sum at most one.  Subtracting the bad pairs and
then all pairs with denominator at most $Q/4$ gives
\begin{align*}
 \#\mathcal R_k(J)
 &\geq \frac{3\lambda Q^2}{32}
 -Q\left(1+\frac\lambda8+\frac14+
                 \left(\frac\lambda2+1\right)k\right)\\
 &\geq\frac{3LQ^2}{64}-3Q(k+1).
 \tag{K7}
\end{align*}
For $k\geq4$, induction gives $2^k\geq k^2$.  Condition (K4) gives
$Lk\geq512$, so
\[
 LQ\geq Lk^2\geq256(k+1)\geq192(k+1).
\]
The last term in (K7) is therefore at most $LQ^2/64$, proving the lower
bound in (K3).

Distinct reduced fractions with denominators at most $Q$ are separated
by at least $Q^{-2}$: their difference has a nonzero integer numerator
over a denominator at most $Q^2$.  In an interval of length $L/2$ this
gives at most $LQ^2/2+1$ fractions.  Our bounds imply $LQ^2\geq2$, so
this is at most $LQ^2$, proving the upper bound in (K3).

\subsection{Small approximation intervals and their intersections}

For $k\geq K(J,N)$ put
\[
 t_k=\min\{2^k\psi(2^k),1\},\qquad
 r_k=\frac{t_k}{8\,4^k},
\]
and let $E_k(J)$ be the union of the intervals of radius $r_k$ centered
at the fractions in $\mathcal R_k(J)$.  An interval with radius zero is
understood as empty.  All these intervals lie in $J$, since their centers
lie in $J^*$ and $r_k\leq1/(8Q^2)\leq L/4$.  They are pairwise disjoint
at each fixed scale, by the separation just proved.  Moreover, for their
labels $p/q$,
\[
 r_k\leq\frac{\psi(Q)}{8Q}\leq\frac{\psi(q)}q,
\]
where the last inequality uses $q\leq Q$ and monotonicity.  Thus every
$E_k(J)$ consists of valid approximations with denominators $q>N$.
Using (K3) and $2r_k=t_k/(4Q^2)$ yields
\[
 \frac{Lt_k}{128}\leq |E_k(J)|\leq\frac{Lt_k}{4}.
 \tag{K8}
\]

Here is the complete overlap estimate.  Fix $k<\ell$ and one interval
$I$ of $E_k(J)$.  A radius-$r_\ell$ interval meeting $I$ has its center
in an interval of length $|I|+2r_\ell$.  Separation at scale $\ell$
bounds the number of such centers by
$(|I|+2r_\ell)4^\ell+1$.  Multiplying by $2r_\ell$ gives
\[
 |I\cap E_\ell(J)|
 \leq\frac{t_\ell}{4}|I|
       +\left(\frac{t_\ell}{4}+1\right)2r_\ell
 \leq\frac{t_\ell}{4}|I|+\frac54(2r_\ell).
\]
Sum this over the disjoint intervals of $E_k(J)$, use (K8), and use
$\#\mathcal R_k(J)\leq L4^k$.  The result is
\[
 |E_k(J)\cap E_\ell(J)|
 \leq\frac L{16}t_kt_\ell
       +\frac{5L}{16}4^{k-\ell}t_\ell.
 \tag{K9}
\]
This is the sole correlation estimate needed.

\subsection{A finite block covers a fixed fraction of every interval}

Take any finite block $k=K,\ldots,B$ of admissible scales and write
$S=\sum_{k=K}^B t_k$.  On the finite partition by all interval endpoints,
let $F=\sum_{k=K}^B\mathbf1_{E_k(J)}$.  The integral signs in what follows
are abbreviations for finite rational sums of segment lengths times
integer values.  Equations (K8)--(K9) give
\[
 \int F\geq LS/128
\]
and
\begin{align*}
 \int F^2
 &=\sum_k|E_k|+2\sum_{k<\ell}|E_k\cap E_\ell|\\
 &\leq L\left(\frac S4+\frac{S^2}{16}+\frac{5S}{24}\right)
  =L\left(\frac{S^2}{16}+\frac{11S}{24}\right).
 \tag{K10}
\end{align*}
We used only $2\sum_{k<\ell}t_kt_\ell\leq S^2$ and the finite geometric
sum $\sum_{k<\ell}4^{k-\ell}\leq1/3$.  When $S\geq1$, the right side
of (K10) is at most $LS^2$.

For completeness, finite weighted Cauchy--Schwarz says
\[
 (\int F)^2\leq |\{F>0\}|\int F^2.
\]
On the endpoint partition this is the elementary identity that the
difference between the right and left sides is
$\sum_{i<j}w_iw_j(v_i-v_j)^2\geq0$, after omitting the cells with value
zero.  Therefore, if $S\geq1$,
\[
 \left|\bigcup_{k=K}^B E_k(J)\right|\geq\frac{L}{16384}.
 \tag{K11}
\]
All divisions have explicitly positive rational denominators.

\subsection{One call to the divergence modulus supplies the block}

Let $K=K(J,N)$, let $s=2^K$, and put
\[
 B=\lceil\log_2D(s)\rceil.
\]
The dyadic blocks $[2^k,2^{k+1}-1]$, $K\leq k\leq B$, include every
integer from $s$ to $D(s)$.  Monotonicity gives
\[
 1\leq\sum_{q=s}^{D(s)}\psi(q)
 \leq\sum_{k=K}^B 2^k\psi(2^k).
\]
For any finite nonnegative rational list $(a_k)$,
\[
 \sum_k\min(a_k,1)\geq\min\left(\sum_k a_k,1\right).
\]
This follows by induction from its two-term version, or by a bounded
case distinction according as some $a_k\geq1$.  Hence $S\geq1$, as
required for (K11).  The cap at one and the factor eight in the radii
are useful: one unit of tail mass suffices, even if concentrated at one
dyadic scale.

\subsection{Amplification by the finitely many residual gaps}

Set $b_0=16384$.  Start with the empty union $U_0$.  At each stage, write
$[0,1]\setminus U_j$ as a finite list of closed rational gaps.  For each
gap of positive length, apply the block construction to its interior
$J$ and add its block union to $U_j$.  Leave singleton gaps in the
complement list.  Freeze the old gap list at the beginning of the round:
process each old gap once, aggregate the new intervals, and only then
compute the next gap list.  All added intervals are inside the gap in which they
were constructed.  If $R_j=|[0,1]\setminus U_j|$, (K11) gives
\[
 R_{j+1}\leq(1-1/b_0)R_j.
\]
The rational binomial inequality
$\bigl(b_0/(b_0-1)\bigr)^{b_0}\geq1+b_0/(b_0-1)>2$
implies $(1-1/b_0)^{b_0}<1/2$.  After
\[
 T=b_0(m+1)
\]
stages we therefore have $R_T\leq2^{-(m+1)}<2^{-m}$.  The maximum of
the finitely many denominator bounds $2^B$ is $M(N,m)$.  This proves
(K2), with a fully specified finite algorithm.

Every loop has an explicit finite bound.  At a scale $k$, enumerate
$1\leq q\leq2^k$ and $0\leq p\leq q$, test coprimality and interval
membership by rational comparisons, and retain the required fractions.
The number of stages is prescribed in advance.  The finite endpoint
sort, gap construction, block endpoint, and new denominator maximum are
primitive recursive in the finite data.  The bounds are intentionally
enormous; efficiency is not the claim.

For membership-sensitive uses, open components that merely touch must
not be merged across their common endpoint.  Merge strict overlaps,
and retain every closed residual gap, including singleton gaps.  Their
zero length may be ignored only in the length calculation.

\subsection{What ``almost every'' means constructively}

The intrinsic iPRA result is (K1) or (K2), expressed in codes for finite
rational data.  Ordinary Lebesgue measure is an external interpretation:
(K1) makes the limsup null in the convergent case, and (K2) makes every
tail union conull, hence their countable intersection conull in the
divergent case.  No countable additivity, density theorem, or zero--one
law was used to prove the finite bounds.

There is also a useful constructive null-test presentation of the
divergent conclusion.  Fix $m$, and for each $n\geq0$ construct
$U(n+1,m+n+3)$.  Its finite closed residual-gap list has total length
less than $\eta_n=2^{-(m+n+3)}$.  If the list has $r_n\geq1$ entries,
enlarge every gap $(a,b)$, including singleton gaps, to
\[
 (a-\delta_n,b+\delta_n),\qquad
 \delta_n=\eta_n/(2r_n).
\]
The total length of this finite list is less than $2\eta_n$.  Concatenate
these lists, or use a primitive recursive pairing of the two list
indices with empty entries for invalid inner indices.  This produces a
primitive recursive open cover $V_m$ whose every finite sublist has
length sum bounded by
\[
 \sum_{n\geq0}2\eta_n=2^{-(m+1)}<2^{-m};
\]
the displayed infinite sum is shorthand for its geometric finite-prefix
bound.

If a represented real avoids $V_m$ for some $m$, it has, for every $N$,
a rational approximation with denominator greater than $N$.  Indeed the
positive rational enlargement margin separates it from every residual
gap of the relevant finite union.  A sufficiently precise rational
approximation to the real therefore locates it inside that union, and
finite interval searching supplies one of the labelled approximation
intervals.  This uses only finite rational case distinctions.  For a
completely explicit extraction bound, subdivide by the finitely many
endpoints, compute the least positive gap between distinct endpoints,
and use a real-name accuracy smaller than one quarter of that gap and
one quarter of the enlargement margin; finite local cases then choose
an original interval with positive slack.  Zero-width endpoint cells
inside the union are covered by an interval crossing the endpoint;
those outside the union belong to the enlarged residual list.

One must not intuitionistically identify the negation of
$\forall N\,\exists q\geq N$ with $\exists N\,\forall q\geq N$.
The finite-coverage statement and the null-test avoidance implication
make no such inference.  Nor do we infer a primitive recursive modulus
from bare convergence or divergence: the relevant modulus is supplied
as part of the constructive input.

\section{Real-valued inputs, positive witnesses, and examples}
\label{sec:extensions}

\subsection{Real-valued approximation functions with rational output covers}

The rational-valued restriction in the finite metric statement is convenient
for exact endpoint calculations.  It can be removed at the represented-input
level without adding an internal sort of real numbers.  Suppose a fixed
primitive recursive function $a(q,s)\in\Q$ names a nonnegative value
$\psi(q)$ with $|a(q,s)-\psi(q)|\le2^{-s}$.  Assume monotonicity of these
named values and a primitive recursive unit-tail modulus $D$ as before.
Every instance of these assumptions can be expressed using inequalities
between rational names.

At scale $Q=2^k$, set
\[
 e_k=2^{-(2k+4)},\qquad
 b_k=\max\{0,a(2^k,2k+4)-e_k\}.
\]
Then
\[
 0\le b_k\le\psi(2^k),\qquad
 0\le\psi(2^k)-b_k\le2e_k.
\]
Use $t_k=\min(2^k b_k,1)$ and the same rational radii
$r_k=t_k/(8\,4^k)$.  These intervals are contained in the original
approximation intervals because $b_k\le\psi(2^k)\le\psi(q)$ for
$q\le2^k$.  Monotonicity of $b_k$ itself is not required.

For a starting scale $K$, put $s=2^K$, $u=D(s)$ and $v=D(u+1)$.
The two disjoint blocks $s,\ldots,u$ and $u+1,\ldots,v$ each have
mass at least one, so the total mass through $v$ is at least two.
For $B=\lceil\log_2 v\rceil$, dyadic condensation yields
\[
 \sum_{k=K}^B2^k\psi(2^k)\ge2.
\]
The loss caused by the rational lower approximations is bounded by
\[
 \sum_{k=K}^B2^k\bigl(\psi(2^k)-b_k\bigr)
 \le\sum_{k=K}^B2^{-k-3}\le\frac14.
\]
Thus $\sum2^k b_k\ge7/4$, and capping at one still gives
$\sum t_k\ge1$.  Every subsequent counting, overlap and amplification
argument is unchanged.  The output is a rational list of intervals,
even when the original target radii are irrational.

For convergence, suppose the named values have a modulus $C(s)\ge1$
with all finite tail sums bounded by $2^{-s}$.  For a requested $m$ use
the rational upper envelope
\[
 u_m(q)=\max\{0,a(q,q+m+4)+2^{-(q+m+4)}\}.
\]
It satisfies $\psi(q)\le u_m(q)\le\psi(q)+2^{-q-m-3}$.
For $q\ge C(m+3)$, the open intervals of radius $u_m(q)/q$
centered at $p/q$, $0\le p\le q$, cover every original approximation.
Their unrestricted lengths, over any finite set of these denominators,
sum to at most
\[
 4\left(2^{-m-3}+\sum_{q\ge1}2^{-q-m-3}\right)
 =2^{-m}.
\]
Again the infinite geometric notation denotes a bound on every finite
prefix.  No positive lower bound for every $\psi(q)$ is needed in either
extension.

\subsection{An explicit extraction lemma for padded finite covers}

There is a constructive point in interpreting the null tests that deserves
a complete finite argument.  The negation of $a<x<b$ does not generally
give a decidable choice between $x\le a$ and $x\ge b$.  We avoid such
choices about real numbers altogether.

Let $U$ be a finite rational open interval list contained in $(0,1)$.
Compute its rational complement cells and the closed residual gaps by
sorting endpoints; include the singleton endpoints $0,1$ if necessary.
Pad every gap on both sides by a fixed positive rational $\delta$, and
write $V$ for that padded finite open cover.  Let $h>0$ be the least
positive difference between distinct endpoints of $U$ together with
$0,1$, and put
\[
 \varepsilon=\frac1{32}\min(\delta,h).
\]
The list contains $0,1$, so $h$ exists and is computed by a finite minimum.
Given a rational approximation $y$ to $x\in[0,1]$ with
$|x-y|<\varepsilon$, the following entirely rational case distinction
finds either an original interval containing $x$ with positive slack,
or a padded gap containing $x$ with positive slack.

First check whether $y$ is within $4\varepsilon$ of an endpoint $c$.
There is at most one distinct such endpoint because distinct endpoints
are at least $32\varepsilon$ apart.  If $c$ lies in an original open
interval, that interval's endpoints are at least $h$ away from $c$;
it therefore contains $x$, whose distance from $c$ is below
$5\varepsilon<h$.  If no original interval contains $c$, then $c$
belongs to a residual gap; its padded interval contains $x$ because
$|x-c|<5\varepsilon<\delta$.  These membership checks are comparisons
of rational endpoints.

If there is no such nearby endpoint, $y$ belongs to the interior of a
single endpoint cell and is more than $4\varepsilon$ from its boundary.
It cannot be outside $[0,1]$, since it is within $\varepsilon$ of $x$.
Then $x$ belongs to the same cell.  The cell is either covered by an
original interval, which supplies the desired interval, or is contained
in a residual gap, whose padding covers $x$.  Its coverage status is
read from the cell's rational midpoint.

Consequently, given the avoidance assumption that $x$ belongs to none of
the padded intervals, the second outcome is impossible and the finite
algorithm returns an original labelled interval containing $x$.  A
precision $s$ with $2^{-s}<\varepsilon$ is obtained by a bounded rational
calculation.  No search for a semidecidable real inequality, Markov's
principle, or excluded middle for real order has been used.

Apply this lemma to each residual list in the null-test construction of
Section~\ref{sec:finite-khintchine}.  Given a specific test index $m$ and
the assertion that $x$ avoids every listed interval of $V_m$, we obtain
a positive rational approximation beyond every prescribed denominator.
For fixed primitive recursive names, moduli and this index $m$, the
output witness is primitive recursive in the denominator threshold.
Iterating with the previous denominator as the next threshold produces
an increasing-denominator sequence of reduced approximants.

\paragraph{What is and is not asserted.}
The effective null test is a primitive recursive enumeration of rational
open intervals for each $m$, with a proved bound $2^{-m}$ on the sum of
the lengths of every finite sublist.  The positive conclusion takes an
explicit avoidance index as data.  We do not turn
$\neg(x\in\bigcap_m V_m)$ into $\exists m\,(x\notin V_m)$ by a
classical quantifier-negation rule.  In ordinary measure theory the
finite estimates imply the usual almost-everywhere theorem.  For arbitrary
classical monotone inputs, the same finite argument works with their
finite mass witnesses; that external interpretation does not assert
that those inputs have primitive recursive moduli.

\subsection{Concrete admissible inputs}

\paragraph{Golden ratio and the sharp coefficient.}
The continued-fraction digits $1,1,1,\ldots$ are primitive recursive.
Alternatively the name obtained by rational bisection of $z^2-5$ and
$I(Q)=4Q^2$ supplies the Hurwitz interface.  The optimality estimate has
a simple visible consequence: for $c=2/5$,
\[
 q^2<\frac{2c^2}{1-5c^2}=\frac85.
\]
Thus $|\phi-p/q|<(2/5)q^{-2}$ can hold only with $q=1$.
The exact coefficient $1/\sqrt5$ does give approximants beyond every
denominator, so retaining it in the constructive theorem is substantive.

\paragraph{Divergent harmonic examples.}
For any fixed positive rational $c$, take $\psi(q)=c/q$ and put
$r=\lceil2/c\rceil$.  The explicit elementary modulus
\[
 D(s)=2^r s-1
\]
works: on each of the $r$ blocks
$2^j s\le q<2^{j+1}s$, the sum of $1/q$ is at least $1/2$.
Hence the sum of $c/q$ from $s$ to $D(s)$ is at least $cr/2\ge1$.
In particular $c=1/4$ gives $D(s)=256s-1$ and target approximation
$|x-p/q|<1/(4q^2)$.  The finite coverage theorem applies to this
genuinely shrinking target, not only to large radii.

\paragraph{A convergent example.}
For $\psi(q)=q^{-2}$, the comparison
\[
 \frac1{q^2}\le\frac1{q(q-1)}=\frac1{q-1}-\frac1q\quad(q\ge2)
\]
gives a finite tail bound $1/(N-1)$.  Thus $C(m)=2^m+1$ is a valid
convergence modulus, and the target approximation is $|x-p/q|<q^{-3}$.
The null-cover conclusion follows with no monotonicity assumption beyond
what this particular example happens to satisfy.

\section{iPRA adequacy and a clear route to complete Lean formalization}
\label{sec:formalization}

\subsection{The internal statements}

Write $\operatorname{HurCert}(B,p,q,s)$ for the conjunction of the
positive-denominator, coprimality, threshold and strict rational tests in
\eqref{eq:hw-certificate}.  After rational arithmetic has been replaced
by its integer-code operations, this is a quantifier-free predicate.
For each fixed pair of primitive recursive definitions $x_s,I$, the
argument constructs primitive recursive output terms $p(B),q(B),s(B)$.
Its arithmetic conclusion is the iPRA sequent
\[
 \operatorname{Name}(x)\wedge\operatorname{Apart}(x,I)
 \ \longrightarrow\
 \forall B\;\operatorname{HurCert}(B,p(B),q(B),s(B)).
\]
Here $\operatorname{Name}$ and $\operatorname{Apart}$ are ordinary
first-order universal formulas about the fixed rational-name functions;
the symbols $x,I$ in this display name those definitions and are not
variables of a new function sort.  The monotonic running maximum for $I$
is part of the primitive recursive construction.

For each fixed rational $\psi,D$, let
$\operatorname{Mon}(\psi)$ express nonnegative monotonic values and let
$\operatorname{Div}(\psi,D)$ be (K-D).  The metric construction supplies
primitive recursive terms $U(N,m),M(N,m)$ and the sequent
\[
 \operatorname{Mon}(\psi)\wedge\operatorname{Div}(\psi,D)
 \ \longrightarrow\
 \forall N,m\,[N\ge1\to\operatorname{Cover}(N,m,U(N,m),M(N,m))].
\]
The predicate $\operatorname{Cover}$ tests that the interval list is
correctly labelled by reduced fractions, every label has $N<q\le M$,
each radius fits the target radius, and its rational residual length is
less than $2^{-m}$.  Each finite-list universal condition is a Boolean
fold, so the final predicate is quantifier-free.  The convergence
sequent similarly asserts the finite length bound under (K-C).

One can first state these as unconditional finite-transcript theorems:
all name evaluations, modulus evaluations and algebraic checks used by a
run are supplied as a finite table; if the table passes the requisite
tests, the output passes the stated certificate test.  Instantiating
the table entries with fixed primitive recursive definitions gives the
schemata above.  This separates arithmetic correctness from any
assumption about the source of the data.

Every code operation must be total on natural numbers.  On malformed
fractions or lists, use a designated default value and a validity bit.
The proofs show that such defaults are never used under the hypotheses.
The reference Python code raises exceptions on invalid inputs for
readability; a literal compilation into PR syntax instead uses these
totalized operations.

\subsection{The induction audit}

The following table gives the finite invariants needed to implement the
human proof.  All comparisons are integer comparisons after clearing
known positive denominators.  Bounded quantifiers in the middle column
are to be compiled into Boolean folds as in Section~\ref{sec:foundations}.

\begin{longtable}{p{0.20\textwidth}p{0.47\textwidth}p{0.23\textwidth}}
\toprule
Computation & Invariant & Induction index\\
\midrule
\endhead
Integer and fraction routines & Sign/denominator validity; cross-product
identities; bounded division remainder and gcd witnesses & Explicit
arithmetic loop counter\\
Convergents & Determinant $\pm1$; positive denominators; growth; stored
rational enclosure has the indicated finite prefix & Prefix length\\
Digit extraction & Fractional-linear denominator bounded away from zero;
image enclosure avoids every relevant integer; endpoint floors agree &
Prefix length and bounded comparison list\\
Hurwitz selection & Chosen pair has $5C^2<1$; precision has its prescribed
margin; selected rational error satisfies the weighted bound & Finite
candidate scan; algebraic identities\\
Counting rational centers & The accumulated count equals the number of
accepted pairs; common-divisor indicators dominate the bad-pair bit &
Bounds on $p,q,d$\\
Dyadic estimates & Running harmonic block sum and telescoping
inverse-square sum have the displayed bounds; $2^k\ge k^2$ for $k\ge4$ &
Integer bound and dyadic exponent\\
Endpoint partition & Sorted endpoints preserve the multiset; each cell
has the correct coverage and incidence bits & Sort and finite cell index\\
Second moment & Accumulated first and second moments satisfy the finite
incidence identities; sum of squares is nonnegative & Cell and pair indices\\
Mass block & Finite grouping of the original tail into dyadic blocks;
capped partial sum dominates the capped total & Dyadic block index\\
Gap filling & Old gaps are processed once; added intervals lie inside
their parent gap and have correct labels; total residual length contracts &
Frozen gap list length, then fixed round count\\
Null covers & Padded list length obeys its geometric budget; finite
rational locator chooses a good interval or a forbidden padded interval &
Finite list and cover enumeration index\\
\bottomrule
\end{longtable}

This table is not a claim that an arbitrary algorithm with a bounded
loop automatically has an iPRA proof.  The algebraic identities,
indicator inequalities, and preservation arguments proving these
particular invariants are given in the preceding sections.  Their
quantifier-free forms are precisely the instances of induction needed.
The finite sums, products and list transformations are primitive
recursive by their displayed recurrences and bounded enumerations.
The large number $16384(m+1)$ of outer rounds is prescribed before the
run.  A round is a fixed primitive recursive transformation of a finite
state; iterating it that many times is again primitive recursion.
New gaps are not processed by an uncontrolled recursive call in the
same round.

The real-number shorthand in the Hurwitz proof has a finite rational
interface, as shown at the end of its extraction argument.  A formal
implementation should prove that interface first and use name coherence
only to supply its enclosures.  It should not build a general classical
real-number theory inside the object language.

\subsection{What was verified in the repository}

The inspected repository snapshot is
\begin{center}
\texttt{34c4c35edc8e648e57587c6faf3a5cfdd7879777}.
\end{center}
Its toolchain is Lean \texttt{v4.32.0}; the mathlib manifest resolves
\texttt{v4.32.0} to the following revision:
\begin{center}
\texttt{81a5d257c8e410db227a6665ed08f64fea08e997}.
\end{center}
The relevant namespace is \texttt{LeanProofs.ConciseIPRA}.  Its source
directory is
\begin{quote}\small
\path{Logic/PrimitiveRecursiveArithmetic/ConciseIPRA/Lean/ConciseIPRA/}.
\end{quote}
The source inspection and exact pinned URLs are recorded in
\path{sources/repository_audit.txt} and \path{sources/source_manifest.json}.

The useful existing interfaces are concrete:\begin{itemize}
\item \texttt{PR.lean} defines the primitive recursive syntax
\texttt{PR}, term syntax \texttt{PTm}, axiom family \texttt{PRAAx}, and
the actual object-theory derived rule \texttt{PRA.qfInd}.
\item \texttt{TeqReflect.lean} supplies characteristic terms
\texttt{chi}, the reflection result \texttt{chi\_iff}, and
\texttt{Teq.qfInd}.  Its quantified formula cases have dummy
characteristic values, so it must not be applied directly to bounded
quantified formulas.  The bounded-fold compiler described above is needed.
\item \texttt{TeqPair.lean} and \texttt{TeqRec.lean} provide pairing and
primitive-recursion machinery, including \texttt{Tm.prc},
\texttt{Q.prc\_zero} and \texttt{Q.prc\_succ}.
\item \texttt{BiInterpretation.lean} proves actual proof translations
between \texttt{TeqAx} and \texttt{PRAAx}.
\item \texttt{NativeBiInterpretation.lean} provides
\texttt{pra\_iff\_native}, translating a \texttt{PRAAx} derivation to
the concise native presentation.
\end{itemize}

The last bridge has the following source-level shape, with its printed
notation expanded informally:
\begin{lstlisting}
(Gamma proves phi in PRAAx) <->
  NPrf (Gamma.map NFm.ofPRA) (NFm.ofPRA phi)
\end{lstlisting}
It can be applied after an actual object-theory proof has been built;
it does not turn an arbitrary ambient Lean theorem into such a proof.
The repository README reports completed foundations without additional
axioms or \texttt{sorry}, with metatheoretic dependencies including
\texttt{propext}, \texttt{Quot.sound}, and, in part of the native chain,
\texttt{Classical.choice}.  Those are dependencies of the Lean
metatheory; the embedded iPRA derivations still use their explicitly
defined intuitionistic rules.  We inspected the relevant source and did
not replay repository builds in this environment.

The \path{docs/drafts/} listing contains a Dirichlet/Kronecker archive,
but its binary contents were not inspected, so no theorem from it is
assumed.  An inspected ordinary Lean finite simultaneous-Dirichlet
module gives a useful pigeonhole pattern; it uses a classical block and
is not itself an embedded iPRA derivation.  Neither main proof in this
report depends on that module.

\subsection{Formalization milestones and completion criteria}

\paragraph{Milestone 1: executable finite arithmetic.}
Implement signed integer and fraction codes, finite list operations,
bounded search, gcd, sorting and interval content.  Prove the primitive
recursive closure certificates as well as their semantic correctness.
The Python artifact gives an executable reference for endpoint behavior
and certificate values; the Lean code should have independent proofs.

\paragraph{Milestone 2: reference Lean theorems.}
Prove the denominator polynomial lemma and weighted-error selector over
$\Z$ and $\Q$.  Prove the finite center-count, overlap and second-moment
lemmas using finite sets or lists and exact rational sums.  Polynomial
normalization tactics can discharge the displayed identities; all
inequalities must be proved from the explicit sign hypotheses.  Build
the rational-enclosure version of the continued-fraction extractor.
Do not use an imported real Hurwitz or metric approximation theorem to
replace these finite lemmas.

\paragraph{Milestone 3: the constructive algorithms.}
Formalize the fixed-round metric state transformation, its local
coverage invariant and contraction.  Formalize the name/modulus
interfaces and the padded-cover locator.  At this point an ambient Lean
theorem proves each finite specification.  This is a substantial useful
checkpoint, but is still separate from the iPRA certification.

\paragraph{Milestone 4: actual iPRA derivations.}
Compile each fixed arithmetic definition to \texttt{PR}/\texttt{PTm}.
Prove bounded-fold reflection inside \texttt{PRAAx}; then implement
the invariants in the table using \texttt{PRA.qfInd}, substitution,
equality and intuitionistic natural deduction.  The final deliverable
for this milestone is a Lean term whose type is an actual
\texttt{Prf PRAAx} sequent for the two displayed schemata (or a
metatheorem generating that term from the fixed primitive recursive
definitions and their hypotheses).  An abstract record assuming these
derivations would only name the remaining task.

\paragraph{Milestone 5: concise native presentation.}
Apply \texttt{pra\_iff\_native} to the constructed derivation.  This uses
the already implemented bridge and introduces no additional number
theory.  Optional classical-comparison modules may interpret the finite
bounds as Lebesgue statements over $\mathbb R$; keep their dependencies
separate from the iPRA result.

The included \path{lean/Blueprint.lean} records finite interfaces and
explicit obligations.  \path{lean/FORMALIZATION_PLAN.txt} describes the
implementation order and repository integration points.  These files
do not claim completion of any of the above milestones beyond their
stated definitions.  On a machine with the pinned repository environment,
the appropriate first check is
\begin{lstlisting}
lake env lean /absolute/path/to/lean/Blueprint.lean
\end{lstlisting}
followed by the separate module builds as the obligations are discharged.

\subsection{Executable checks supplied with this report}

The command \texttt{python3 code/verify.py} regenerates the exact
\path{code/verification.json} record.  Its checks include 9,830 denominator
triples; six Hurwitz certificates for the golden ratio and $\sqrt2$;
1,000 seeded union comparisons against an independent endpoint-cell
calculation; 1,000 intersection comparisons; 1,000 cap inequalities;
90 local overlap estimates; and explicit endpoint edge cases.  All
calculations use integers and exact fractions.  These computations
check implementations and examples; they do not replace the general
proofs.

For instance the golden-ratio example with threshold $B=10$ returns
$p/q=233/144$ and a precision-$39$ rational certificate whose left side is
\[
 \frac{18446391806604141645}{18446744073709551616}<1.
\]
This tests a case close to the sharp constant.  The finite reduced union
for $\psi(q)=1/q$ and $5\le q\le80$ has 1,960 labelled intervals and
exact covered length
\[
 \frac{55336826409}{59204622400}.
\]
Neither this sample denominator cutoff nor the moderate local tests are
claimed to instantiate the astronomical worst-case cutoff in the proof.
The full general residual-gap algorithm and the general modulus-based
continued-fraction extractor are specified in the manuscript; they have
not been executed end to end by the sample program.

\clearpage
\begin{thebibliography}{9}
\bibitem{hurwitz}
A. Hurwitz, \emph{Ueber die angen\"aherte Darstellung der Irrationalzahlen
durch rationale Br\"uche}, Mathematische Annalen \textbf{39} (1891),
279--284. \url{https://doi.org/10.1007/BF01206656}.
\bibitem{filaseta}
M. Filaseta, \emph{Math 784: Algebraic Number Theory, Instructor's Notes},
Theorem 25, printed pp.~20--21.
\url{https://people.math.sc.edu/filaseta/gradcourses/TheMath784Notes.pdf}.
\bibitem{bv}
V. Beresnevich and S. Velani, \emph{Ubiquity and a general logarithm law
for geodesics}, arXiv:0707.1225 (2007).
\url{https://arxiv.org/abs/0707.1225}.
\bibitem{bdv}
V. Beresnevich, D. Dickinson and S. Velani,
\emph{Measure theoretic laws for lim sup sets}, arXiv:math/0401118,
version 3 (2004). \url{https://arxiv.org/abs/math/0401118}.
\bibitem{repo}
V. Reshetnikov and contributors, \emph{ProveIt: Concise presentations of
intuitionistic primitive recursive arithmetic}. Repository snapshot
\texttt{34c4c35edc8e648e57587c6faf3a5cfdd7879777}; source inspection
15 November 2026. Exact pinned source links and toolchain versions are
included in \path{sources/source_manifest.json}.
\url{https://github.com/VladimirReshetnikov/ProveIt/tree/main/Logic/PrimitiveRecursiveArithmetic/ConciseIPRA}.
\bibitem{requested}
Requested statement references: \emph{Hurwitz's theorem (number theory)}
and the Khinchin section of \emph{Diophantine approximation}, Wikipedia,
accessed 15 November 2026. The proofs above are self-contained;
these links identify the requested classical formulations.
\url{https://en.wikipedia.org/wiki/Hurwitz\%27s_theorem_(number_theory)};
\url{https://en.wikipedia.org/wiki/Diophantine_approximation}.
\end{thebibliography}

\end{document}
